%%%%%%%%%%%%%%%%%%%%%%%%%%%%%%%%%%%%%%%%%%%%%%%%%%%%%%%%%%%%%
\section{模型检验与结果复核}

本章依据各问模型类型选择必要检验，不机械同时安排误差分析和灵敏度分析。

% 优化问题：逐条复核约束、边界、重叠或资源限制；比较上下界、精确小规模解、
% 多初值或邻近不可行搜索，并明确“未找到”不等于“证明不可行”。
% 预测/分类：外样本切分、残差、多指标、类别不平衡和基线。
% 评价：一致性、权重扰动和排序稳定性。
% 机理：量纲、边界条件、参数辨识和外部数据对照。

\subsection{结果正确性与约束检验}

针对……，重新计算……，并逐项检查……约束。检验结果见表……。
结果表明……；该检验能够证明……，但不能证明……。

\subsection{稳定性、误差或灵敏度分析}

选择对结论最有解释意义的……进行……扰动／多初值／外样本检验。
扰动范围、评价指标与结果见表……。在……范围内，……变化为……，
因此关于……的结论稳健／仅在……条件下成立。

% 若该类检验不适用，删除本节并说明采用了何种替代验证，不为凑章节生成空表。
